
\documentclass[10pt]{beamer}
\usepackage[utf8]{inputenc}
\usepackage[T5]{fontenc}
\usepackage[vietnamese]{babel}
\usepackage{lmodern}
\usepackage{graphicx}
\usepackage{booktabs}
\usepackage{array}
\usepackage{xcolor}
\usepackage{tikz}
\usetikzlibrary{arrows.meta,positioning,fit,shapes.geometric,shapes.multipart,calc}
\usepackage{hyperref}
\AtBeginDocument{%
  \DeclareFontShape{T5}{lmss}{m}{it}{%
    <-8.5> t5-lmsso8
    <8.5-9.5> t5-lmsso9
    <9.5-11> t5-lmsso10
    <11-15.5> t5-lmsso12
    <15.5-> t5-lmsso17
  }{}
}

% Chọn theme
\usetheme{Madrid}
\usecolortheme{default}
\setbeamertemplate{navigation symbols}{}
\setbeamertemplate{footline}[frame number]

\definecolor{prismblue}{RGB}{42,91,155}
\definecolor{prismgreen}{RGB}{50,130,90}
\definecolor{prismorange}{RGB}{200,120,30}

\newcommand{\projectname}{GocTriThuc}
\newcommand{\maybeimage}[2]{%
  \IfFileExists{#1}{\includegraphics[width=#2\textwidth]{#1}}{%
  \begin{tikzpicture}[scale=#2]
    \draw[rounded corners=5pt,thick,fill=gray!8] (0,0) rectangle (12,6.2);
    \draw[fill=prismblue!15,draw=prismblue] (0.35,5.55) rectangle (11.65,5.95);
    \draw[fill=white,draw=gray!50] (0.65,0.55) rectangle (11.35,5.25);
    \node[align=center,text width=8cm] at (6,3) {Ảnh minh họa\\\texttt{#1}};
  \end{tikzpicture}}
}

% Thông tin trang bìa
\title[Nền tảng GocTriThuc]{\texorpdfstring{BÁO CÁO BÀI TẬP LỚN \\ NHẬP MÔN CÔNG NGHỆ PHẦN MỀM}{BÁO CÁO BÀI TẬP LỚN - NHẬP MÔN CÔNG NGHỆ PHẦN MỀM}}
\subtitle{Đề tài: Nền tảng học tập trực tuyến GocTriThuc}
\author[Nhóm GocTriThuc]{\texorpdfstring{%
  \textbf{Giảng viên hướng dẫn:} TS. Nguyễn Quốc Tuấn \\ \vspace{0.25cm}
  \textbf{Nhóm thực hiện:} \\
  Trần Tuấn Anh (202416124) \and Nguyễn Công Vinh (202400083) \\
  Phạm Văn Sâm (202400114) \and Vũ Hoàng Tuấn (202400080) \\
  Lê Thành Trung (202400076)
}{Nhóm GocTriThuc}}
\institute[HUST]{\texorpdfstring{Đại học Bách Khoa Hà Nội \\ Trường CNTT \& TT}{Đại học Bách Khoa Hà Nội - Trường CNTT & TT}}
\date{Ngày 16 tháng 6 năm 2026}

\begin{document}

% Slide 1: Bìa
\begin{frame}
    \titlepage
\end{frame}

% Slide 2: Mục lục
\begin{frame}{Nội dung trình bày}
    \tableofcontents
\end{frame}

% --- PHẦN 1: TỔNG QUAN ---
\section{Tổng quan dự án}

\begin{frame}{Lý do chọn đề tài \& Mục tiêu}
    \begin{block}{Lý do chọn đề tài}
        \begin{itemize}
            \item Các nền tảng hiện tại như Moodle, Canvas cấu hình phức tạp, khó tùy biến cho nhóm nhỏ.
            \item Cần một hệ thống e-learning tinh gọn, dễ sử dụng, có kiến trúc rõ ràng và kiểm soát dữ liệu nội bộ.
        \end{itemize}
    \end{block}
    \begin{block}{Mục tiêu dự án}
        \begin{itemize}
            \item \textbf{Giao diện:} thân thiện, responsive, xây dựng bằng React/Vite.
            \item \textbf{API:} backend RESTful chuẩn hóa bằng Spring Boot.
            \item \textbf{Dữ liệu \& bảo mật:} PostgreSQL, OAuth2/JWT, phân quyền theo role/permission.
            \item \textbf{Quy trình:} phát triển Agile/Scrum, đóng gói Docker.
        \end{itemize}
    \end{block}
\end{frame}

\begin{frame}{Phân tích hiện trạng \& Giải pháp tương tự}
    \begin{table}[]
        \centering
        \resizebox{\textwidth}{!}{%
        \begin{tabular}{@{}p{3cm}p{4.5cm}p{4.5cm}@{}}
            \toprule
            \textbf{Nền tảng} & \textbf{Ưu điểm} & \textbf{Nhược điểm / Hạn chế} \\ \midrule
            \textbf{Google Classroom} & Dễ sử dụng, ít chi phí. & Khó tùy biến workflow, báo cáo học tập còn hạn chế. \\ \midrule
            \textbf{Moodle / Canvas} & Đầy đủ tính năng, hệ sinh thái lớn. & Cấu hình nặng, giao diện phức tạp, chi phí triển khai và vận hành cao. \\ \midrule
            \textbf{GocTriThuc} \newline (Dự án) & \textbf{Kiến trúc hiện đại, tinh gọn, tùy biến cao cho các nhóm/lớp học nhỏ.} & Đang trong giai đoạn phát triển, chưa có app Mobile native. \\ \bottomrule
        \end{tabular}%
        }
    \end{table}
    \vspace{0.3cm}
    $\Rightarrow$ \textit{Dự án hướng tới sự cân bằng: \textbf{Dễ dùng như Google Classroom nhưng có hệ thống khóa học/bài test chuyên sâu như Moodle}.}
\end{frame}

\begin{frame}{Phương pháp luận phát triển}
    \begin{block}{Áp dụng mô hình Agile/Scrum}
        \begin{itemize}
            \item \textbf{Lý do:} Phù hợp với hệ thống E-learning cần thay đổi linh hoạt theo phản hồi của người dùng. Giảm rủi ro so với mô hình Waterfall.
            \item \textbf{Tổ chức công việc:} Phân chia thành các Sprint ngắn (1-2 tuần).
            \item \textbf{Quy trình mỗi Sprint:}
            \begin{itemize}
                \item Lập kế hoạch (Sprint Planning) \& Phân rã User Story.
                \item Lập trình Frontend/Backend song song (thiết kế API Contract).
                \item Review mã nguồn, kiểm thử tích hợp.
                \item Bàn giao chức năng chạy được (Definition of Done).
            \end{itemize}
        \end{itemize}
    \end{block}
\end{frame}

% --- PHẦN 2: PHÂN TÍCH YÊU CẦU ---
\section{Đặc tả yêu cầu hệ thống}

\begin{frame}{Đối tượng người dùng (Actors)}
    Hệ thống phục vụ 4 nhóm người dùng chính:
    \vspace{0.2cm}
    \begin{itemize}
        \item \textbf{Khách vãng lai:} xem thông tin, khóa học công khai, đăng nhập Google OAuth.
        \item \textbf{Học viên:} ghi danh, học Video/Blog, làm Test, xem tiến độ, bình luận.
        \item \textbf{Giảng viên:} tạo/quản lý khóa học, module, bài học, ngân hàng câu hỏi, duyệt yêu cầu.
        \item \textbf{Quản trị viên:} quản lý người dùng, phân quyền, giám sát hệ thống.
    \end{itemize}
\end{frame}

\begin{frame}{Yêu cầu chức năng cốt lõi}
    \begin{itemize}
        \item \textbf{Xác thực \& hồ sơ:} đăng nhập Google OAuth, quản lý profile.
        \item \textbf{Quản lý khóa học:} tạo/xuất bản khóa học, module, bài học Video/Blog/Test.
        \item \textbf{Tương tác học tập:} ghi danh Public/Restricted, theo dõi tiến độ.
        \item \textbf{Kiểm tra đánh giá:} ngân hàng câu hỏi, tạo bài test, lưu phiên làm bài.
        \item \textbf{Trao đổi:} đăng thông báo, bình luận phân cấp.
        \item \textbf{Dashboard:} phân hệ quản lý riêng cho Giảng viên và Admin.
    \end{itemize}
\end{frame}

\begin{frame}{Use Case Diagram tổng thể}
\centering
\resizebox{!}{0.78\textheight}{%
\begin{tikzpicture}[
    actor/.style={
        draw,
        ellipse,
        fill=orange!10,
        minimum width=2.4cm,
        align=center,
        font=\small
    }, 
    uc/.style={
        draw,
        ellipse,
        fill=blue!6,
        minimum width=3.8cm,
        minimum height=0.8cm,
        align=center,
        font=\small
    }, 
    every node/.style={font=\small}
]

\node[actor] (guest) at (-2.5, 4.85) {Khách};
\node[actor] (student) at (-2.5, 1.0) {Học viên};
\node[actor] (admin) at (-2.5, -5.5) {Admin};
\node[actor] (teacher) at (10.5, -4.25) {Giảng viên};

\node[draw, thick, rounded corners, minimum width=9cm, minimum height=16cm, label={[anchor=south, font=\large]north:Hệ Thống}] (boundary) at (4, -1.65) {};

\node[uc] (login) at (5, 5.5) {Đăng nhập OAuth};
\node[uc] (catalog) at (5, 4.2) {Xem danh mục};
\node[uc] (enroll) at (5, 2.9) {Ghi danh khóa học};
\node[uc] (learn) at (5, 1.6) {Học lesson};
\node[uc] (test) at (5, 0.3) {Làm bài test};
\node[uc] (comment) at (5, -1.0) {Bình luận};
\node[uc] (profile) at (5, -2.3) {Cập nhật hồ sơ};
\node[uc] (course) at (5, -3.6) {Quản lý khóa học};
\node[uc] (lesson) at (5, -4.9) {Quản lý module/lesson};
\node[uc] (question) at (5, -6.2) {Ngân hàng câu hỏi};
\node[uc] (announce) at (5, -7.5) {Đăng thông báo};
\node[uc] (useradmin) at (5, -8.8) {Quản trị người dùng};

\foreach \a/\b in {guest/login, guest/catalog, student/catalog, student/enroll, student/learn, student/test, admin/useradmin, teacher/course, teacher/lesson, teacher/question}{
    \draw (\a) -- (\b);
}

\draw (student) to[bend right=15] (comment.west);
\draw (student) to[bend right=25] (profile.west);
\draw (admin) to[bend left=25] (profile.west);
\draw (teacher) to[bend right=30] (comment.east);
\draw (teacher) to[bend left=30] (announce.east);

\draw[dashed, -{Latex}] (test) -- node[pos=0.5, right, font=\scriptsize] {<<extend>>} (learn);
\draw[dashed, -{Latex}] (course) -- node[pos=0.5, right, font=\scriptsize] {<<include>>} (lesson);
\draw[dashed, -{Latex}] (course.west) .. controls ++(-2,0) and ++(-2,0) .. node[left, font=\scriptsize] {<<include>>} (question.west);
\draw[dashed, -{Latex}] (course.west) .. controls ++(-2.5,0) and ++(-2.5,0) .. node[left, font=\scriptsize] {<<include>>} (announce.west);
\end{tikzpicture}}
\end{frame}

% --- PHẦN 3: THIẾT KẾ HỆ THỐNG ---
\section{Thiết kế hệ thống}

\begin{frame}{Kiến trúc tổng thể}
    \begin{center}
    \resizebox{0.96\textwidth}{!}{%
    \begin{tikzpicture}[
        node distance=1.5cm, 
        box/.style={draw,rounded corners,align=center,minimum width=3.2cm,minimum height=1cm,fill=blue!6}, 
        db/.style={draw,cylinder,shape border rotate=90,aspect=0.25,align=center,minimum height=1.3cm,fill=green!8}
    ]

    \node[box] (browser) {Trình duyệt\\React + Vite};
    \node[box, right=3.2cm of browser] (api) {REST API\\Spring Boot Controllers};
    \node[box, right=1cm of api] (service) {Service Layer\\Business Rules};
    \node[db, right=2.5cm of service] (db) {PostgreSQL};

    \node[box, below=1.5cm of api] (oauth) {Google OAuth2\\Identity Provider};
    \node[box, below=1.5cm of service] (file) {File Storage\\Avatar/Resources};

    \draw[-{Latex}] (browser) -- node[above=1mm] {HTTPS/JSON} (api);
    \draw[-{Latex}] (api) -- (service);
    \draw[-{Latex}] (service) -- node[above=1mm] {JPA/SQL} (db);
    \draw[-{Latex}] (api) -- (oauth);
    \draw[-{Latex}] (service) -- (file);
    \end{tikzpicture}}
    \end{center}
    \begin{itemize}
        \item \textbf{Frontend:} React + Vite, giao tiếp qua HTTPS/JSON.
        \item \textbf{Backend:} Spring Boot, REST API, Service Layer.
        \item \textbf{Database:} PostgreSQL, Flyway migrations, JPA/SQL.
    \end{itemize}
\end{frame}

\begin{frame}{Công nghệ sử dụng \& Triển khai}
    \begin{columns}[T]
        \begin{column}{0.5\textwidth}
            \begin{block}{Frontend}
                \begin{itemize}
                    \item \textbf{Framework:} ReactJS (Vite)
                    \item \textbf{Ngôn ngữ:} TypeScript
                    \item \textbf{Styling:} Tailwind CSS / UI Components
                \end{itemize}
            \end{block}
            \begin{block}{Backend}
                \begin{itemize}
                    \item \textbf{Framework:} Spring Boot 3.x
                    \item \textbf{Ngôn ngữ:} Java
                    \item \textbf{Bảo mật:} Spring Security, OAuth2, JWT
                \end{itemize}
            \end{block}
        \end{column}
        
        \begin{column}{0.5\textwidth}
            \begin{block}{Cơ sở dữ liệu \& Lưu trữ}
                \begin{itemize}
                    \item \textbf{Database:} PostgreSQL
                    \item \textbf{Migration:} Flyway
                    \item \textbf{File Storage:} Tích hợp Cloud / Local
                \end{itemize}
            \end{block}
            \begin{block}{Môi trường \& Quy trình}
                \begin{itemize}
                    \item \textbf{Đóng gói:} Docker, Docker Compose
                    \item \textbf{Tài liệu API:} Swagger / OpenAPI
                    \item \textbf{Quản lý mã nguồn:} Git / GitHub
                \end{itemize}
            \end{block}
        \end{column}
    \end{columns}
\end{frame}

\begin{frame}{Thiết kế Cơ sở dữ liệu}
\centering
\resizebox{0.98\textwidth}{!}{%
\begin{tikzpicture}[entity/.style={draw,rounded corners,fill=green!7,minimum width=2.8cm,minimum height=0.8cm,align=center}, node distance=1.2cm and 1.6cm, every node/.style={font=\small}]
\node[entity] (users) {users};
\node[entity,right=of users] (courses) {courses};
\node[entity,right=of courses] (modules) {modules};
\node[entity,right=of modules] (lessons) {lessons};
\node[entity,below=of users] (roles) {roles / permissions};
\node[entity,below=of courses] (enroll) {enrollments};
\node[entity,below=of modules] (ann) {announcements};
\node[entity,below=of lessons] (comments) {comments};
\node[entity,above=of lessons] (tests) {lesson\_tests};
\node[entity,above=of modules] (questions) {questions};
\node[entity,above=of courses] (sessions) {test\_sessions};
\draw[-{Latex}] (courses) -- node[above]{1-n} (modules);
\draw[-{Latex}] (modules) -- node[above]{1-n} (lessons);
\draw[-{Latex}] (users) -- node[above]{owner} (courses);
\draw[-{Latex}] (users) -- (roles);
\draw[-{Latex}] (users) -- (enroll);
\draw[-{Latex}] (courses) -- (enroll);
\draw[-{Latex}] (courses) -- (ann);
\draw[-{Latex}] (lessons) -- (comments);
\draw[{Latex}-{Latex}] (lessons) -- (tests);
\draw[-{Latex}] (users) -- (sessions);
\draw[-{Latex}] (tests.north) -- ++(0,0.8) -| (sessions.north);
\draw[-{Latex}] (tests) -- (questions);
\end{tikzpicture}}
\vspace{1mm}
\begin{itemize}\scriptsize
    \item Sơ đồ ERD rút gọn sử dụng trực tiếp cấu trúc từ báo cáo chính.
    \item Trung tâm gồm \texttt{users}, \texttt{courses}, \texttt{modules}, \texttt{lessons}, \texttt{lesson\_tests} và \texttt{test\_sessions}.
\end{itemize}
\end{frame}

\begin{frame}{Class Diagram rút gọn}
\centering
\resizebox{0.98\textwidth}{!}{%
\begin{tikzpicture}[class/.style={draw,rectangle split,rectangle split parts=3,minimum width=3.6cm,align=left,fill=gray!5}, node distance=1.2cm and 1.5cm, every node/.style={font=\scriptsize}]
\node[class] (user) {\textbf{User}\nodepart{second}id: Long\\email: String\\displayName: String\\avatarUrl: String\nodepart{third}+getRoles()\\+updateProfile()};
\node[class,right=of user] (course) {\textbf{Course}\nodepart{second}id: Long\\title: String\\visibility: Enum\\ownerId: Long\nodepart{third}+publish()\\+updateInfo()};
\node[class,right=of course] (module) {\textbf{Module}\nodepart{second}id: Long\\courseId: Long\\title: String\\order: int\nodepart{third}+reorder()};
\node[class,below=of module] (lesson) {\textbf{Lesson}\nodepart{second}id: Long\\moduleId: Long\\type: LessonType\\order: int\nodepart{third}+markComplete()};
\node[class,below=of course] (test) {\textbf{LessonTest}\nodepart{second}id: Long\\statement: String\\timeLimit: int\\settings: JSON\nodepart{third}+startSession()\\+grade()};
\node[class,below=of user] (question) {\textbf{Question}\nodepart{second}id: Long\\content: String\\type: Enum\nodepart{third}+validateAnswer()};
\node[class,below=of lesson] (comment) {\textbf{Comment}\nodepart{second}id: Long\\content: String\\parentId: Long\nodepart{third}+reply()\\+edit()};
\draw[-{Latex}] (course) -- node[above]{owner} (user);
\draw[-{Latex}] (module) -- node[above]{belongs to} (course);
\draw[-{Latex}] (lesson) -- node[right]{belongs to} (module);
\draw[-{Latex}] (test) -- node[above]{extends} (lesson);
\draw[-{Latex}] (test) -- node[above]{uses} (question);
\draw[-{Latex}] (comment) -- node[right]{on} (lesson);
\end{tikzpicture}}
\end{frame}

\begin{frame}{Sơ đồ tuần tự: Học viên làm bài kiểm tra}
\centering
\resizebox{0.98\textwidth}{!}{%
\begin{tikzpicture}[lifeline/.style={draw,minimum width=2.5cm,minimum height=0.8cm,fill=blue!6}, msg/.style={-{Latex}}, every node/.style={font=\small}]
\node[lifeline] (stu) at (0,0) {Học viên};
\node[lifeline] (fe) at (3.5,0) {React UI};
\node[lifeline] (ctrl) at (7,0) {TestSessionController};
\node[lifeline] (svc) at (10.8,0) {TestService};
\node[lifeline] (db) at (14.2,0) {PostgreSQL};
\foreach \x in {0,3.5,7,10.8,14.2}{\draw[dashed] (\x,-0.5) -- (\x,-7.2);} 
\draw[msg] (0,-1) -- node[above]{Bấm bắt đầu} (3.5,-1);
\draw[msg] (3.5,-1.6) -- node[above]{POST /test-sessions} (7,-1.6);
\draw[msg] (7,-2.2) -- node[above]{startSession()} (10.8,-2.2);
\draw[msg] (10.8,-2.8) -- node[above]{insert session} (14.2,-2.8);
\draw[msg] (14.2,-3.4) -- node[above]{sessionId} (10.8,-3.4);
\draw[msg] (10.8,-4.0) -- node[above]{TestSessionResponse} (7,-4.0);
\draw[msg] (7,-4.6) -- node[above]{JSON} (3.5,-4.6);
\draw[msg] (0,-5.2) -- node[above]{Chọn đáp án} (3.5,-5.2);
\draw[msg] (3.5,-5.8) -- node[above]{PUT /answers} (7,-5.8);
\draw[msg] (7,-6.4) -- node[above]{saveAnswer()} (10.8,-6.4);
\draw[msg] (10.8,-7.0) -- node[above]{upsert answer} (14.2,-7.0);
\end{tikzpicture}}
\end{frame}

% --- PHẦN 4: GIAO DIỆN & KIỂM THỬ ---
\section{Giao diện \& Kiểm thử}


\begin{frame}{Giao diện: Đăng nhập}
\centering
\maybeimage{screen1.png}{0.66}\\[1mm]
\scriptsize Đăng nhập / trang vào hệ thống
\end{frame}

\begin{frame}{Giao diện: Danh sách khóa học}
\centering
\maybeimage{screen2.png}{0.66}\\[1mm]
\scriptsize Danh sách hoặc trang khám phá khóa học
\end{frame}

\begin{frame}{Giao diện: Chi tiết khóa học}
\centering
\maybeimage{screen3.png}{0.66}\\[1mm]
\scriptsize Chi tiết khóa học / classroom
\end{frame}

\begin{frame}{Giao diện: Xem bài học}
\centering
\maybeimage{screen4.png}{0.66}\\[1mm]
\scriptsize Xem bài học, tài nguyên và bình luận
\end{frame}

\begin{frame}{Giao diện: Dashboard giảng viên}
\centering
\maybeimage{screen5.png}{0.66}\\[1mm]
\scriptsize Dashboard giảng viên / công cụ quản lý nội dung
\end{frame}

\begin{frame}{Giao diện: Dashboard quản trị viên}
\centering
\maybeimage{screen6.png}{0.66}\\[1mm]
\scriptsize Dashboard quản trị viên / kiểm tra và quản trị
\end{frame}

\begin{frame}{Chiến lược Kiểm thử}
    \centering
    \resizebox{\textwidth}{!}{%
    \begin{tabular}{@{}lll@{}}
        \toprule
        \textbf{Mức kiểm thử} & \textbf{Phạm vi} & \textbf{Kỹ thuật \& Tiêu chí} \\ \midrule
        \textbf{Unit} & Service, Util, Mapper & Mock repo, test nhánh nghiệp vụ. \\ 
        \textbf{Integration} & Controller, Repository, DB & Gọi API, kiểm tra HTTP status, JSON. \\ 
        \textbf{E2E} & Luồng nghiệp vụ chính & Khách $\rightarrow$ Ghi danh $\rightarrow$ Học $\rightarrow$ Test. \\ 
        \textbf{Security} & Endpoint phân quyền & Test token sai/thiếu, check lỗi 401/403. \\ 
        \textbf{Hiệu năng} & Tải danh sách, submit test & Phản hồi trong ngưỡng cho phép. \\ \bottomrule
    \end{tabular}}
\end{frame}

% --- PHẦN 5: TỔNG KẾT ---
\section{Kết luận \& Hướng phát triển}

\begin{frame}{Kết luận \& Hướng phát triển}
    \begin{block}{Kết quả đạt được}
        \begin{itemize}
            \item Hoàn thiện đặc tả và thiết kế hệ thống e-learning với các tính năng cốt lõi.
            \item Xây dựng kiến trúc hiện đại: React, Spring Boot, PostgreSQL, Docker/CI-CD.
            \item Thiết kế database bám sát Flyway migrations và mô hình phân quyền rõ ràng.
        \end{itemize}
    \end{block}
    \begin{block}{Hướng phát triển tương lai}
        \begin{itemize}
            \item Chứng chỉ, huy hiệu học tập và gamification.
            \item AI chấm tự luận, gợi ý lộ trình học tập, analytics nâng cao.
            \item Ứng dụng di động hoặc PWA, audit log và bảo mật production.
        \end{itemize}
    \end{block}
\end{frame}

\begin{frame}
    \begin{center}
        \Huge \textbf{CẢM ƠN THẦY VÀ CÁC \\BẠN ĐÃ LẮNG NGHE!} \\
        \vspace{1cm}
        \large Q\&A
    \end{center}
\end{frame}

\end{document}
